\chapter{Apr.~2 --- Applications of Riemann-Roch}

\section{Applications of Riemann-Roch}

\begin{example}
  Let $C = \PP^1$. We have
  $g(\PP^1) = h^1(\PP^1, \OO_{\PP^1}) = 0$.
  Using the isomorphism
  \begin{align*}
    \Cl(\PP^1) &\overset{\cong}{\longrightarrow}
    \Z \\
    D &\longmapsto \deg D,
  \end{align*}
  we have $\OO_{\PP^1}(D) = \OO_{\PP^1}(d)$
  with $d = \deg D$. So
  \[
    \chi(\PP^1, \OO_{\PP^1}(D))
    = \chi(\PP^1, \OO_{\PP^1}(d))
    = h^0(\PP^1, \OO_{\PP^1}(d))
    - h^1(\PP^1, \OO_{\PP^1}(d))
    = d + 1 = \deg D + 1 - g.
  \]
\end{example}

\begin{example}
  Let $C \subseteq \PP^2$ be a smooth
  plane curve of degree $d$. Then
  we have seen that
  \[
    g = h^1(C, \OO_C) = \frac{(d - 1)(d - 2)}{2}.
  \]
\end{example}

\begin{corollary}\label{cor:linearly_equivalent_divisors_have_same_degree}
  If $D_1$ and $D_2$ are linearly
  equivalent divisors on $C$, then
  $\deg D_1 = \deg D_2$.
\end{corollary}

\begin{proof}
  If $D_1 \sim D_2$, then
  $\OO_C(D_1) \cong \OO_C(D_2)$, so
  $\chi(C, \OO_C(D_1)) = \chi(C, \OO_C(D_2))$,
  which then implies that
  $\deg D_1 = \deg D_2$ by the
  Riemann-Roch theorem.
\end{proof}

\begin{definition}
  Let $\mathcal{L}$ be an invertible sheaf
  on $C$. The \emph{degree} of $\mathcal{L}$
  is $\deg \mathcal{L} := \deg D$ where
  $D$ is a divisor such that
  $\mathcal{L} \cong \OO_C(D)$.
\end{definition}

\begin{remark}
  The degree of $\mathcal{L}$ is
  well-defined by Corollary
  \ref{cor:linearly_equivalent_divisors_have_same_degree}.
\end{remark}

\begin{example}
  Let $\mathcal{L}$ be an invertible sheaf
  on $C$. If $h^0(C, \mathcal{L}) \ge 1$,
  then there exists $0 \ne s \in H^0(C, \mathcal{L})$.
  Thus there exists $D \ge 0$ such that
  $\mathcal{L} \cong \OO_C(D)$, so
  $\deg \mathcal{L} = \deg D \ge 0$.
  Furthermore, if $h^0(C, \mathcal{L}) \ge 1$
  and $\deg \mathcal{L} = 0$, then
  $\mathcal{L} \cong \OO_C$.
\end{example}

\begin{corollary}\label{cor:degree_of_canonical_divisor}
  $\deg \omega_C = 2 g - 2$.
\end{corollary}

\begin{proof}
  By Serre duality, we have
  $\chi(C, \OO_C) = -\chi(C, \omega_C)$
  (as $h^0(\OO_C) = h^1(\omega_C)$ and
  $h^1(\OO_C) = h^0(\omega_C)$). So by
  the Riemann-Roch theorem, we have
  \[
    -(1 - g)
    = -\chi(C, \OO_C)
    = \chi(C, \omega_C)
    = \deg \omega_C + 1 - g,
  \]
  so $\deg \omega_C = 2g - 2$, as desired.
\end{proof}

\begin{remark}
  If $0 \ne \eta \in \Gamma(C, \omega_C)$, then
  $\omega_C \cong \OO_C(\divv_C \eta)$, so
  by Corollary \ref{cor:degree_of_canonical_divisor},
  \[
    2g - 2
    = \deg \divv_C \eta = \text{\# zeros of $\eta$}.
  \]
\end{remark}

\begin{definition}
  We say that a
  divisor $D$ on $C$ is \emph{non-special} if
  $h^0(C, \OO_C(K - D)) = 0$
  (or equivalently, by Serre duality, if
  $h^1(C, \OO_C(D)) = 0$).
\end{definition}

\begin{remark}
  For $D$ non-special, we have
  $h^0(C, \OO_C(D)) = \deg D + 1 - g$ by
  Riemann-Roch.
\end{remark}

\begin{lemma}
  If $\deg D \ge 2g - 1$, then
  $D$ is non-special.
\end{lemma}

\begin{proof}
  If $\deg D \ge 2g - 1$, then
  $\deg(K - D) \le (2g - 2) - (2g - 1) = -1$,
  so $h^0(C, \OO_C(K - D)) = 0$.
\end{proof}

\begin{remark}
  If $\deg D \ge 2g - 1$, then
  $D$ is non-special, which implies
  $h^0(C, \OO_C(D)) = \deg D + 1 - g$.
\end{remark}

\section{More on Embeddings to Projective Space}

\begin{prop}\label{prop:globally_generated_iff_h0_drops_by_1}
  An invertible sheaf
  $\mathcal{L}$ on $C$ is globally
  generated if and only if
  \[
    h^0(C, \mathcal{L})
    - h^0(C, \mathcal{L}(-P)) = 1
  \]
  for every $P \in C$.
\end{prop}

\begin{remark}
  We define $\mathcal{L}(-P) := \mathcal{L} \otimes \OO_C(-P)$.
  In these cases, the fibre will be denoted by
  \[
    \mathcal{L}_{(P)}
    := \mathcal{L}_P / \mathfrak{m}_P \mathcal{L}_P
    = H^0(C, \mathcal{L} / \mathcal{I}_P \mathcal{L}).
  \]
\end{remark}

\begin{proof}[Proof of Proposition \ref{prop:globally_generated_iff_h0_drops_by_1}]
  For $P \in C$, we have a short exact sequence
  \[
    \begin{tikzcd}
      0 \ar[r] & \OO_C(-P)
      \ar[r] & \OO_C \ar[r] & \OO_P \ar[r] & 0
    \end{tikzcd}
  \]
  Applying $- \otimes \mathcal{L}$, we
  get an exact sequence (since
  $\mathcal{L}$ is invertible)
  \[
    \begin{tikzcd}
      0 \ar[r] & \mathcal{L}(-P)
      \ar[r] & \mathcal{L} \ar[r] & \mathcal{L} \otimes \OO_P
      \ar[r] & 0.
    \end{tikzcd}
  \]
  Taking $H^0$, we then get an exact sequence
  \[
    \begin{tikzcd}[row sep=small]
      0 \ar[r] & H^0(\mathcal{L}(-P)) \ar[r] & H^0(\mathcal{L})
      \ar[r, "\phi"] & H^0(\mathcal{L} \otimes \OO_P) \\
             & & & \mathcal{L}_{(P)} \cong k \ar[u, equal]
    \end{tikzcd}
  \]
  So $\mathcal{L}$ is globally generated
  at $P$
  if and only if $\phi$ is surjective,
  if and only if
  $h^0(\mathcal{L}) - h^0(\mathcal{L}(-P)) = 1$.
\end{proof}

\begin{prop}\label{prop:very_ample_iff_h0_drops_by_2}
  An invertible sheaf $\mathcal{L}$ on $C$
  is very ample if and only if
  \[
    h^0(C, \mathcal{L})
    - h^0(C, \mathcal{L}(-P - Q)) = 2
  \]
  for every $P, Q \in C$ (not necessarily distinct).
\end{prop}

\begin{proof}
  By the proof of Proposition
  \ref{prop:globally_generated_iff_h0_drops_by_1},
  we have
  \[
    h^0(C, \mathcal{L})
    - h^0(C, \mathcal{L}(-P))
    = 1
  \]
  for all $P \in C$, so
  $\mathcal{L}$ is globally generated.
  In this setting, we know that
  $\mathcal{L}$ is very ample if and only if
  $\mathcal{L}$ separates points and
  tangent vectors. For $P \ne Q \in C$,
  consider the short exact sequence
  \[
    \begin{tikzcd}
      0 \ar[r] & \OO_C(-P - Q)
      \ar[r] & \OO_C
      \ar[r] & \OO_P \oplus \OO_Q
      \ar[r] & 0
    \end{tikzcd}
  \]
  Applying $- \otimes \mathcal{L}$
  and $H^0$, we get an exact sequence
  \[
    \begin{tikzcd}
      0 \ar[r] & H^0(\mathcal{L}(-P - Q)) \ar[r] & H^0(\mathcal{L}) \ar[r, "\psi"]
               & \mathcal{L}_{(P)} \oplus \mathcal{L}_{(Q)}
    \end{tikzcd}
  \]
  So $\mathcal{L}$ separates $P, Q$ if and only
  if
  $\psi$ is surjective, if and only if
  $h^0(C, \mathcal{L})
    - h^0(C, \mathcal{L}(-P - Q)) = 2$.

  Next fix $P \in C$, and consider the
  short exact sequence
  \[
    \begin{tikzcd}
      0 \ar[r] & \OO_C(-2P)
      \ar[r] & \OO_C
      \ar[r] & \OO_C / \OO_C(-2P) \ar[r] & 0
    \end{tikzcd}
  \]
  Again applying
  $- \otimes \mathcal{L}$ and $H^0$, we
  get an exact sequence
  \[
    \begin{tikzcd}
      0 \ar[r] & H^0(\mathcal{L}(-2P))
      \ar[r] & H^0(\mathcal{L})
      \ar[r, "\gamma"] & \mathcal{L}_P / \m_P^2 \mathcal{L}_P
    \end{tikzcd}
  \]
  (note that
  $\OO_C / \OO_C(-2P) = i_P(\OO_{C, P} / \m_P^2)$).
  Thus $\mathcal{L}$ separates tangent
  vectors at $P$ if and only if
  $\gamma$ is surjective, if and only if
  $h^0(C, \mathcal{L})
    - h^0(C, \mathcal{L}(-2P)) = 2$.
\end{proof}

\begin{corollary}\label{cor:globally_generated_and_very_ample_if_degree_large_enough}
  We have the following:
  \begin{enumerate}
    \item if $\deg \mathcal{L} \ge 2g$, then
      $\mathcal{L}$ is globally generated;
    \item if $\deg \mathcal{L} \ge 2g + 1$,
      then $\mathcal{L}$ is very ample.
  \end{enumerate}
\end{corollary}

\begin{proof}
  (1) If $\deg \mathcal{L} \ge 2g$, then
  \[
    h^0(C, \mathcal{L})
    = \deg \mathcal{L} + 1 - g
    \quad \text{and} \quad
    h^0(C, \mathcal{L}(-P))
    = \deg \mathcal{L}(-P) + 1 - g
  \]
  as $\mathcal{L}$ and $\mathcal{L}(-P)$ are
  both non-special. So
  $h^0(C, \mathcal{L})
    - h^0(C, \mathcal{L}(-P)) = 1$ for all
  $P \in C$, so $\mathcal{L}$ is globally generated
  by Proposition \ref{prop:globally_generated_iff_h0_drops_by_1}.

  (2) The proof is similar, use that
  $\deg \mathcal{L} \ge 2g + 1$ implies
  \[
    \deg \mathcal{L}
    > \deg \mathcal{L}(-P)
    \ge \deg \mathcal{L}(-P - Q)
    \ge 2g - 1,
  \]
  so $\mathcal{L}, \mathcal{L}(-P), \mathcal{L}(-P - Q)$ are all
  non-special. So
  \[
    h^0(C, \mathcal{L})
    - h^0(C, \mathcal{L}(-P - Q))
    = \deg \mathcal{L} - \deg \mathcal{L}(-P - Q)
    = 2,
  \]
  so $\mathcal{L}$ is very ample
  by Proposition \ref{prop:very_ample_iff_h0_drops_by_2}.
\end{proof}

\begin{remark}
  Corollary \ref{cor:globally_generated_and_very_ample_if_degree_large_enough}
  is sharp: Let $C = \PP^1$, then
  $\deg \OO_{\PP^1} = 2g$, but
  $\OO_{\PP^1}$ is not very ample. Similarly,
  $\deg \OO_{\PP^1}(-1) = 2g - 1$, but
  $\OO_{\PP^1}(-1)$ is not globally generated.
\end{remark}

\section{Applications to Curves}

\begin{corollary}
  An invertible sheaf $\mathcal{L}$ on $C$ is
  ample if and only if $\deg \mathcal{L} > 0$.
\end{corollary}

\begin{proof}
  $(\Rightarrow)$ If $\mathcal{L}$ is ample,
  then there exists $m > 0$ such that
  $\mathcal{L}^{\otimes m}$ is very ample, so
  \[
    m \deg \mathcal{L}
    = \deg \mathcal{L}^{\otimes m} > 0,
  \]
  which implies that $\deg \mathcal{L} > 0$.

  $(\Leftarrow)$ If $\mathcal{L} > 0$, then
  $\deg \mathcal{L}^{\otimes (2g + 1)} = (\deg \mathcal{L})(2g + 1) \ge 2g + 1$,
  so $\mathcal{L}$ is very ample.
\end{proof}

\begin{corollary}
  If the genus of $C$ is $g = 0$, then
  $C \cong \PP^1$.
\end{corollary}

\begin{proof}
  Fix $P \in C$ and set $\mathcal{L} = \OO_C(P)$.
  Since $\deg L = 1 \ge g + 1$, we know
  $\mathcal{L}$ is very ample. Furthermore,
  \[
    h^0(\mathcal{L})
    = \deg \mathcal{L} + 1 - g
    = 1 + 1 - 0 = 2.
  \]
  So we get an embedding
  $C \hookrightarrow \PP^{2 - 1} = \PP^1$.
  Since $\dim C = \dim \PP^1 = 1$ and
  $\PP^1$ is irreducible,
  $C \cong \PP^1$.
\end{proof}

\begin{corollary}
  If the genus of $C$ is $g = 1$, then
  $C$ is isomorphic to a cubic curve in $\PP^2$.
\end{corollary}

\begin{proof}
  Fix $P \in C$ and set $\mathcal{L} = \OO_C(3P)$.
  Since $\deg \mathcal{L} = 3 \ge 2g + 1$,
  we know
  $\mathcal{L}$ is very ample.
  Similarly,
  \[
    h^0(\mathcal{L})
    = \deg \mathcal{L} + 1 - g = 3.
  \]
  Thus we get an embedding
  $C \hookrightarrow \PP^{3 - 1} = \PP^2$.
  So $C$ is a plane curve of some
  degree $d$. Using that
  \[
    1 = g = \frac{1}{2}(d - 1)(d - 2),
  \]
  we see that we can only have $d = 3$.
\end{proof}
